\appendix

\section{Diagnostic Methods for the Canonical Data Contract}\label{app:canonical-diagnostic-chain}

This appendix organizes the canonical data-contract diagnosis by error source and scientific interpretation.

\begin{table}[!ht]
\centering
\caption{Diagnosis of anomalies and isolation of error sources in the canonical data contract.}
\label{tab:canonical-diagnostic-chain}
\scriptsize
\begin{tabularx}{\textwidth}{@{}l>{\raggedright\arraybackslash}X>{\raggedright\arraybackslash}X@{}}
\toprule
Diagnostic step & Observation and analysis & Scientific interpretation \\
\midrule
Same-trajectory comparison & After removing a cross-trajectory comparison, mode1 met the criterion whereas mixed did not. & The apparent positive result was caused by inconsistent comparison objects. \\
Repeated single-mode case & A repeated mode1 case could also exhibit closure error. & Modal mixing was not necessary for the anomaly. \\
Cotangent projection & Insufficient \(N_{\mathrm{diag}}=2\) diagnostic truncation was not the cause, and projected momentum closed more accurately. & Momentum should be projected to the cotangent space of the fixed basis. \\
Solver-output semantics & The internal generalized-alpha stage was not available through the public API, whereas BDF \texttt{tsteps} returned same-solver-step states. & BDF is used only for state sampling and label generation. \\
Fixed-basis initialization & Applying the normalization/sign transform to displacement and velocity gave \(12/12\) consistent initial conditions and short-window closure for mixed. & The anomaly arose from inconsistent initialization in the fixed basis. \\
Two-level time-step comparison & G2 increased from \(0.00147337\) to \(0.00155680\). & Two levels did not meet the predefined monotonic-convergence criterion. \\
Three levels and \(1\,\mathrm{s}\) validation & Observed orders were \(1.64186\)--\(2.00088\), and the extended-case closure and energy quantities met their thresholds. & The evidence supports only this single diagnostic sentinel as a candidate input to a later A1 qualification experiment. \\
\bottomrule
\end{tabularx}
\end{table}


The diagnosis successively examines comparison objects, modal mixing, reduced-order truncation, the cotangent-momentum definition, solver-output semantics, and fixed-basis initialization. The false positive caused by a cross-trajectory comparison shows that closure errors must be computed on the same physical trajectory and at the same sampling times. A repeated single-mode case can still exhibit the anomaly, so the error cannot be attributed directly to modal mixing. Increasing the reduced representation did not reveal the signature of an \(N_{\mathrm{diag}}=2\) diagnostic-truncation error; projected momentum closed more accurately than \(M_N\dot q_{\mathrm{native}}\).

Because the internal generalized-alpha stage is unavailable through the public API, exported displacement, velocity, and load cannot be guaranteed to correspond to the same internal solution stage. Order-2 BDF with \texttt{tsteps} is used to obtain states and vector-field labels at the same solver step. This choice does not imply that BDF is symplectic and does not interpret the discrete trajectory as a symplectic map.

Force closure and perturbation-energy drift use independent criteria. An earlier force comparison was sensitive to endpoints and the momentum anchor. On common-support indices \(2\ldots998\), with a fixed \(p[2]\) anchor, the force-v2 joint L2 is below \(2\%\) for all three trajectories. This result does not alter the trajectory-level conclusions obtained with the original \(0.5\%\) Bridge perturbation-energy criterion.

\section{Metric Definitions and Numerical Settings}\label{app:numerical-settings}

\subsection{Canonical mapping and closure quantities}

For the fixed \(N_{\mathrm{diag}}=2\) diagnostic basis, the mapping is
\begin{align}
q &= \Pi_{\mathrm{conf},N}(U), &
p &= B_N(q)^{T}P_{\mathrm{full}}, \\
\dot q &= M_N(q)^{-1}p, &
\dot p &= \dot B_N(q)^{T}P_{\mathrm{full}}+B_N(q)^{T}\dot P_{\mathrm{full}} .
\end{align}
Here \(q\in\mathbb{R}^{N_{\mathrm{diag}}}\) is a diagnostic configuration coordinate and \(p\in T_q^*\mathcal{Q}_{N}\) is its associated cotangent coordinate, with \(B_N,\dot B_N\in\mathbb{R}^{m_f\times N_{\mathrm{diag}}}\) and \(P_{\mathrm{full}},\dot P_{\mathrm{full}}\in\mathbb{R}^{m_f}\). If \(L_{\mathrm{full}}:=\dot P_{\mathrm{full}}\) is introduced, the final line may be written as \(\dot p=\dot B_N^{T}P_{\mathrm{full}}+B_N^{T}L_{\mathrm{full}}\), but the transport term cannot be omitted. Closure errors compare the exported \(\dot q\) with \(M_N^{-1}p\) and the exported \(\dot p\) with this transport-corrected momentum rate, while the time-discretization study also examines \(q\), \(p\), and total energy.

The three-level study uses coarse, middle, and fine grids. If the errors on two adjacent levels are \(e_h\) and \(e_{h/2}\), the observed order is
\[
r_{\mathrm{obs}}=\log_2\!\left(\frac{e_h}{e_{h/2}}\right).
\]
The resulting observed orders span \(1.64186\)--\(2.00088\). In the \(1\,\mathrm{s}\) case, the thresholds for G2, pointwise G3, and impulse G3 are all \(0.02\), and the peak-to-peak perturbation-energy-drift threshold is \(0.005\).

\subsection{Bridge and low-dimensional network experiments}

Bridge perturbation-energy drift is compared trajectory by trajectory with the \(0.5\%\) criterion and the exploratory \(1\%\) criterion. Force-v2 uses an independent \(2\%\) joint-L2 criterion. The energy and force quantities are not combined into a single conclusion.

P0 uses a four-dimensional nonlinear coupled oscillator. The training, validation, and test sets contain 80, 16, and 20 trajectories, respectively; each trajectory has 80 steps, \(\Delta t=0.025\), and the models are trained for 300 epochs. Five random seeds estimate algorithmic repeatability.

P1 and P1.5 each comprise \(60/60\) completed training-and-evaluation runs, and P1.6 comprises \(30/30\) hard-constrained models. The three physical trajectories---mode1, mode2, and mixed12---are partitioned by leave-one-whole-trajectory-out. Different random seeds share the same physical trajectory set and therefore do not increase the physical sample size. P1.5 compares physics-only and residual learning, whereas P1.6 compares predictive accuracy and structural residuals before and after imposing the hard equilibrium constraint.

\subsection{Subsequent numerical instantiation}

Multi-amplitude statistics, the full 96-trajectory data set, and the P2 model comparison will retain whole-trajectory partitioning and report cross-condition means, variances, and paired comparisons. Open-port studies additionally require power-balance, dissipation-sign, and port-leakage quantities, while high-fidelity studies must quantify the discrepancy of the Kirchhoff model relative to Timoshenko-beam and solid models. Numerical KAM analysis requires nondegeneracy and Diophantine margins, high-order errors, symplectic-map defects, and invariant-torus residuals in a common norm. Implementation code and supplementary data will be provided with the accompanying code and supplementary materials.

The current clean checkout contains only the V1--V3 \texttt{verification\_skeleton.py} and its \texttt{verification\_results.json} output. They are not formal CHNN, a formal COMSOL--Galerkin bridge, formal force-v2 code, a P0/P1/P2 production experiment, a controller, an adaptive law, or a numerical KAM certificate. Removing this boundary requires formal source code, configuration, an environment lock, raw arrays, a manifest, clean-run logs, automatically generated tables and figures, and provenance-consistency records; while those materials are missing, A-17 remains PARTIAL.
